
\documentclass{article}
\usepackage{colortbl}
\usepackage{makecell}
\usepackage{multirow}
\usepackage{supertabular}

\begin{document}

\newcounter{utterance}

\centering \large Interaction Transcript for game `dond', experiment `coop\_en', episode 13 with Qwen3.5{-}27B.
\vspace{24pt}

{ \footnotesize  \setcounter{utterance}{1}
\setlength{\tabcolsep}{0pt}
\begin{supertabular}{c@{$\;$}|p{.15\linewidth}@{}p{.15\linewidth}p{.15\linewidth}p{.15\linewidth}p{.15\linewidth}p{.15\linewidth}}
    \# & $\;$A & \multicolumn{4}{c}{Game Master} & $\;\:$B\\
    \hline

    \theutterance \stepcounter{utterance}  
    & & \multicolumn{4}{p{0.6\linewidth}}{
        \cellcolor[rgb]{0.9,0.9,0.9}{
            \makecell[{{p{\linewidth}}}]{
                \texttt{\tiny{[P1$\langle$GM]}}
                \texttt{You are playing a cooperative negotiation game in which you have to agree on how to divide a set of items among you and another player.} \\
\\ 
\texttt{Rules:} \\
\texttt{(a) You and the other player are given a set of items. Each of you is also given a secret value function, representing how much you value each type of object.} \\
\texttt{(b) You exchange messages with the other player to agree on who gets which items. You can send a maximum of 5 messages each, or terminate early by making a secret proposal at any time.} \\
\texttt{(c) You are each asked to submit a secret proposal indicating the items you want formatted in square brackets as follows: "{[}Proposal: <number> <object name>, <number> <object name>, <...>{]}"} \\
\texttt{(d) If your proposals are complementary, i.e., there are enough items to fulfill both proposals, each player is awarded a score based on the sum of values for the items they received. Otherwise, both of you get zero points.} \\
\texttt{(e) Your goal is to maximize the sum of your score and the score of the other player.} \\
\\ 
\texttt{Let us start.} \\
\\ 
\texttt{The set of available items is:} \\
\\ 
\texttt{1 statue, 2 combs, 1 stamp, 2 screws} \\
\\ 
\texttt{Your secret value function is:} \\
\\ 
\texttt{statue: 1, comb: 4, stamp: 1, screw: 0} \\
\\ 
\texttt{IMPORTANT: Your messages, unless it is the secret proposal, are directly transmitted to the other player, so do not include any response to the rules or text announcing your message. To make a secret proposal, use the indicated format. Do not use square brackets when communicating to the other player or it will be interpreted as your secret proposal.} \\
            }
        }
    }
    & \\ \\

    \theutterance \stepcounter{utterance}  
    & \multicolumn{4}{p{0.6\linewidth}}{
        \cellcolor[rgb]{0.9,0.9,0.9}{
            \makecell[{{p{\linewidth}}}]{
                \texttt{\tiny{[P1$\rangle$GM]}}
                \texttt{Hi! Let's figure out the best way to split these items to maximize our combined score.} \\
\\ 
\texttt{I value the combs the most (4 points each), while the statue and stamp are worth very little to me (1 point each), and I don't value the screws at all (0 points).} \\
\\ 
\texttt{Could you share how much you value each item? That way, we can give the items to the person who values them the most. For example, if you value the statue or stamp highly, I'd be happy to let you have them in exchange for the combs.} \\
            }
        }
    }
    & & \\ \\

    \theutterance \stepcounter{utterance}  
    & & & \multicolumn{2}{p{0.3\linewidth}}{
        \cellcolor[rgb]{0.9,0.9,0.9}{
            \makecell[{{p{\linewidth}}}]{
                \texttt{\tiny{[GM$|$GM]}}
                \texttt{continue} \\
            }
        }
    }
    & & \\ \\

    \theutterance \stepcounter{utterance}  
    & & & \multicolumn{2}{p{0.3\linewidth}}{
        \cellcolor[rgb]{0.9,0.9,0.9}{
            \makecell[{{p{\linewidth}}}]{
                \texttt{\tiny{[GM$|$GM]}}
                \texttt{Hi! Let's figure out the best way to split these items to maximize our combined score.} \\
\\ 
\texttt{I value the combs the most (4 points each), while the statue and stamp are worth very little to me (1 point each), and I don't value the screws at all (0 points).} \\
\\ 
\texttt{Could you share how much you value each item? That way, we can give the items to the person who values them the most. For example, if you value the statue or stamp highly, I'd be happy to let you have them in exchange for the combs.} \\
            }
        }
    }
    & & \\ \\

    \theutterance \stepcounter{utterance}  
    & & \multicolumn{4}{p{0.6\linewidth}}{
        \cellcolor[rgb]{0.9,0.9,0.9}{
            \makecell[{{p{\linewidth}}}]{
                \texttt{\tiny{[GM$\rangle$P2]}}
                \texttt{You are playing a cooperative negotiation game in which you have to agree on how to divide a set of items among you and another player.} \\
\\ 
\texttt{Rules:} \\
\texttt{(a) You and the other player are given a set of items. Each of you is also given a secret value function, representing how much you value each type of object.} \\
\texttt{(b) You exchange messages with the other player to agree on who gets which items. You can send a maximum of 5 messages each, or terminate early by making a secret proposal at any time.} \\
\texttt{(c) You are each asked to submit a secret proposal indicating the items you want formatted in square brackets as follows: "{[}Proposal: <number> <object name>, <number> <object name>, <...>{]}"} \\
\texttt{(d) If your proposals are complementary, i.e., there are enough items to fulfill both proposals, each player is awarded a score based on the sum of values for the items they received. Otherwise, both of you get zero points.} \\
\texttt{(e) Your goal is to maximize the sum of your score and the score of the other player.} \\
\\ 
\texttt{Let us start.} \\
\\ 
\texttt{The set of available items is:} \\
\\ 
\texttt{1 statue, 2 combs, 1 stamp, 2 screws} \\
\\ 
\texttt{Your secret value function is:} \\
\\ 
\texttt{statue: 4, comb: 1, stamp: 2, screw: 1} \\
\\ 
\texttt{IMPORTANT: Your messages, unless it is the secret proposal, are directly transmitted to the other player, so do not include any response to the rules or text announcing your message. To make a secret proposal, use the indicated format. Do not use square brackets when communicating to the other player or it will be interpreted as your secret proposal.} \\
\\ 
\\ 
\texttt{Hi! Let's figure out the best way to split these items to maximize our combined score.} \\
\\ 
\texttt{I value the combs the most (4 points each), while the statue and stamp are worth very little to me (1 point each), and I don't value the screws at all (0 points).} \\
\\ 
\texttt{Could you share how much you value each item? That way, we can give the items to the person who values them the most. For example, if you value the statue or stamp highly, I'd be happy to let you have them in exchange for the combs.} \\
            }
        }
    }
    & \\ \\

    \theutterance \stepcounter{utterance}  
    & & & \multicolumn{4}{p{0.6\linewidth}}{
        \cellcolor[rgb]{0.9,0.9,0.9}{
            \makecell[{{p{\linewidth}}}]{
                \texttt{\tiny{[GM$\langle$P2]}}
                \texttt{Thanks for sharing your values! That helps a lot.} \\
\\ 
\texttt{Here is how I value the items:} \\
\texttt{{-} Statue: 4 points} \\
\texttt{{-} Comb: 1 point} \\
\texttt{{-} Stamp: 2 points} \\
\texttt{{-} Screw: 1 point} \\
\\ 
\texttt{Looking at our values together, it seems like we have a great opportunity to maximize our combined score by trading based on our preferences:} \\
\texttt{{-} You value the **combs** highly (4 each), while I only value them at 1.} \\
\texttt{{-} I value the **statue** highly (4), while you only value it at 1.} \\
\texttt{{-} I value the **stamp** at 2, which is higher than your 1.} \\
\texttt{{-} We both value the **screws** at 1, so it doesn't matter much who gets them, but since you don't value them at all, I can take them.} \\
\\ 
\texttt{Based on this, here is a proposed split that maximizes our total points:} \\
\texttt{{-} **You get**: 2 combs (worth 8 to you)} \\
\texttt{{-} **I get**: 1 statue (worth 4 to me), 1 stamp (worth 2 to me), and 2 screws (worth 2 to me)} \\
\\ 
\texttt{This way, you get the items you value most, and I get the items I value most. The total combined score would be 8 (your combs) + 4 (my statue) + 2 (my stamp) + 2 (my screws) = 16 points.} \\
\\ 
\texttt{Does this sound good to you? If so, I can make a secret proposal to finalize this.} \\
            }
        }
    }
     \\ \\

    \theutterance \stepcounter{utterance}  
    & & & \multicolumn{2}{p{0.3\linewidth}}{
        \cellcolor[rgb]{0.9,0.9,0.9}{
            \makecell[{{p{\linewidth}}}]{
                \texttt{\tiny{[GM$|$GM]}}
                \texttt{continue} \\
            }
        }
    }
    & & \\ \\

    \theutterance \stepcounter{utterance}  
    & & & \multicolumn{2}{p{0.3\linewidth}}{
        \cellcolor[rgb]{0.9,0.9,0.9}{
            \makecell[{{p{\linewidth}}}]{
                \texttt{\tiny{[GM$|$GM]}}
                \texttt{Thanks for sharing your values! That helps a lot.} \\
\\ 
\texttt{Here is how I value the items:} \\
\texttt{{-} Statue: 4 points} \\
\texttt{{-} Comb: 1 point} \\
\texttt{{-} Stamp: 2 points} \\
\texttt{{-} Screw: 1 point} \\
\\ 
\texttt{Looking at our values together, it seems like we have a great opportunity to maximize our combined score by trading based on our preferences:} \\
\texttt{{-} You value the **combs** highly (4 each), while I only value them at 1.} \\
\texttt{{-} I value the **statue** highly (4), while you only value it at 1.} \\
\texttt{{-} I value the **stamp** at 2, which is higher than your 1.} \\
\texttt{{-} We both value the **screws** at 1, so it doesn't matter much who gets them, but since you don't value them at all, I can take them.} \\
\\ 
\texttt{Based on this, here is a proposed split that maximizes our total points:} \\
\texttt{{-} **You get**: 2 combs (worth 8 to you)} \\
\texttt{{-} **I get**: 1 statue (worth 4 to me), 1 stamp (worth 2 to me), and 2 screws (worth 2 to me)} \\
\\ 
\texttt{This way, you get the items you value most, and I get the items I value most. The total combined score would be 8 (your combs) + 4 (my statue) + 2 (my stamp) + 2 (my screws) = 16 points.} \\
\\ 
\texttt{Does this sound good to you? If so, I can make a secret proposal to finalize this.} \\
            }
        }
    }
    & & \\ \\

    \theutterance \stepcounter{utterance}  
    & & \multicolumn{4}{p{0.6\linewidth}}{
        \cellcolor[rgb]{0.9,0.9,0.9}{
            \makecell[{{p{\linewidth}}}]{
                \texttt{\tiny{[P1$\langle$GM]}}
                \texttt{Thanks for sharing your values! That helps a lot.} \\
\\ 
\texttt{Here is how I value the items:} \\
\texttt{{-} Statue: 4 points} \\
\texttt{{-} Comb: 1 point} \\
\texttt{{-} Stamp: 2 points} \\
\texttt{{-} Screw: 1 point} \\
\\ 
\texttt{Looking at our values together, it seems like we have a great opportunity to maximize our combined score by trading based on our preferences:} \\
\texttt{{-} You value the **combs** highly (4 each), while I only value them at 1.} \\
\texttt{{-} I value the **statue** highly (4), while you only value it at 1.} \\
\texttt{{-} I value the **stamp** at 2, which is higher than your 1.} \\
\texttt{{-} We both value the **screws** at 1, so it doesn't matter much who gets them, but since you don't value them at all, I can take them.} \\
\\ 
\texttt{Based on this, here is a proposed split that maximizes our total points:} \\
\texttt{{-} **You get**: 2 combs (worth 8 to you)} \\
\texttt{{-} **I get**: 1 statue (worth 4 to me), 1 stamp (worth 2 to me), and 2 screws (worth 2 to me)} \\
\\ 
\texttt{This way, you get the items you value most, and I get the items I value most. The total combined score would be 8 (your combs) + 4 (my statue) + 2 (my stamp) + 2 (my screws) = 16 points.} \\
\\ 
\texttt{Does this sound good to you? If so, I can make a secret proposal to finalize this.} \\
            }
        }
    }
    & \\ \\

    \theutterance \stepcounter{utterance}  
    & \multicolumn{4}{p{0.6\linewidth}}{
        \cellcolor[rgb]{0.9,0.9,0.9}{
            \makecell[{{p{\linewidth}}}]{
                \texttt{\tiny{[P1$\rangle$GM]}}
                \texttt{That sounds perfect! Your analysis is spot on. By giving you the statue, stamp, and screws, and keeping the two combs for myself, we maximize the total value to 16 points.} \\
\\ 
\texttt{I agree to this split. Please go ahead and make the secret proposal to finalize it.} \\
\\ 
\texttt{{[}Proposal: 2 comb{]}} \\
            }
        }
    }
    & & \\ \\

    \theutterance \stepcounter{utterance}  
    & & & \multicolumn{2}{p{0.3\linewidth}}{
        \cellcolor[rgb]{0.9,0.9,0.9}{
            \makecell[{{p{\linewidth}}}]{
                \texttt{\tiny{[GM$|$GM]}}
                \texttt{proposal} \\
            }
        }
    }
    & & \\ \\

    \theutterance \stepcounter{utterance}  
    & & & \multicolumn{2}{p{0.3\linewidth}}{
        \cellcolor[rgb]{0.9,0.9,0.9}{
            \makecell[{{p{\linewidth}}}]{
                \texttt{\tiny{[GM$|$GM]}}
                [0, 2, 0, 0]
            }
        }
    }
    & & \\ \\

    \theutterance \stepcounter{utterance}  
    & & \multicolumn{4}{p{0.6\linewidth}}{
        \cellcolor[rgb]{0.9,0.9,0.9}{
            \makecell[{{p{\linewidth}}}]{
                \texttt{\tiny{[GM$\rangle$P2]}}
                \texttt{The other player has submitted a secret proposal. It is now your turn to submit a proposal of your own.} \\
            }
        }
    }
    & \\ \\

    \theutterance \stepcounter{utterance}  
    & & & \multicolumn{4}{p{0.6\linewidth}}{
        \cellcolor[rgb]{0.9,0.9,0.9}{
            \makecell[{{p{\linewidth}}}]{
                \texttt{\tiny{[GM$\langle$P2]}}
                \texttt{{[}Proposal: 2 comb{]}} \\
            }
        }
    }
     \\ \\

    \theutterance \stepcounter{utterance}  
    & & & \multicolumn{2}{p{0.3\linewidth}}{
        \cellcolor[rgb]{0.9,0.9,0.9}{
            \makecell[{{p{\linewidth}}}]{
                \texttt{\tiny{[GM$|$GM]}}
                \texttt{proposal} \\
            }
        }
    }
    & & \\ \\

    \theutterance \stepcounter{utterance}  
    & & & \multicolumn{2}{p{0.3\linewidth}}{
        \cellcolor[rgb]{0.9,0.9,0.9}{
            \makecell[{{p{\linewidth}}}]{
                \texttt{\tiny{[GM$|$GM]}}
                [0, 2, 0, 0]
            }
        }
    }
    & & \\ \\

    \theutterance \stepcounter{utterance}  
    & & & \multicolumn{2}{p{0.3\linewidth}}{
        \cellcolor[rgb]{0.9,0.9,0.9}{
            \makecell[{{p{\linewidth}}}]{
                \texttt{\tiny{[GM$|$GM]}}
                [[0, 2, 0, 0], [0, 2, 0, 0]]
            }
        }
    }
    & & \\ \\

\end{supertabular}
}

\end{document}
